\chapter[Sermon 46. St. John Measures the Temple, and Shows That God Hath a Care of It: and the Choir He Excommunicateth.]{St. John Measures the Temple, and Shows That God Hath a Care of It: and the Choir He Excommunicateth.}
\label{sermon:046}
\markright{Sermon 46. St. John Measures the Temple, and Shows That God Hath ...}

And a reed like a rod was given to me, and I was told, “Rise and measure the temple of God, and the altar, and those who worship there, and the choir within the temple; cast them out and do not measure them, for they are given to the Gentiles, and the holy city they will trample underfoot for forty-two months.”

\index{Antichrist}\index{Turks}The Lord continues comforting, and describing the hostile war against Antichrist, showing that the church will not be forsaken in those Antichristian and Turkish difficulties. The enemies will never so quietly enjoy all things, but the church will also have her champions or defenders, who will most valiantly resist Christ's adversaries.

\index{Papacy}\index{Daniel (Book of)}\index{Ezekiel (Book of)}\index{Rome}And those things are figurative, which are rehearsed in the beginning of the chapter, and seem to be taken from the fortieth chapter of Ezekiel. As are also those which are spoken of in the seventh chapter concerning the faithful sealed, from the ninth chapter of the same prophet. For he is commanded to measure the temple and to cast out the inner sanctuary; of which he shows the cause. And he does not mean the Temple of Jerusalem, which lay in ruins, nor should it be repaired after the prophecy of Daniel and Christ, but the very church of God, that is, the whole number of the elect. For Paul calls the faithful the Temple of God, living indeed, as also Peter in 1 Peter 2 and 1 Corinthians 3 and 2 Corinthians 6. We have said often that Christ is the only altar in the church, and sacrifice for sin, and priest and intercessor on the right hand of the Father. The worshippers are those who worship God through Christ in spirit and truth, and serve him lawfully or with fear. So many as are such, that is, whoever cleaves to Christ, the only peacemaker of the faithful, and serves God truly by faith, they are the very Temple of God and the true church. John has measured these, so that we should understand how the Lord sets his mind to build up the church, not to destroy it. For those who build measure the foundation, upon which the buildings should be set, as appears in Ezekiel 40. Then the temple was also destroyed by the Chaldeans, as the church is now devastated by the Papists and Turks. But the Lord promises by this measuring that he will repair the ruins of the church through the merit of Christ and faithful worshippers; moreover, he signifies that the faithful in these troubles are numbered (as we heard they were sealed) and sure, whom no hostile power can hurt in all these difficulties. For as the altar, Christ, is undefiled and cannot be polluted or destroyed by any power of the Devil, so are the sheep of Christ known to God and do not perish, as also the same Lord Jesus Christ testifies in John 10 and the apostle in 2 Timothy 2. Briefly, the faithful of Christ are in communion with God and with all his good things, in his care, building, number, and defense. This is a most assured consolation. However, where the Lord in the gospel prophesied that the true faithful would be excommunicated by false teachers, he also foresaw what would happen to the ungodly pastors of the false bishops. He says that they do not belong to the building of God, but are to be excommunicated by God, so that the godly should not fear their censure and cursing. And here is the distinction of two sorts, the first of which is more approved: namely, the hall or sanctuary that is within, cast out; that is to say, declare those who are in this sanctuary to be cast out from God. Indeed, the Antichristians will be within the temple, or the inner parts of the temple, and the chiefest part of the church, so that whoever does not acknowledge them and does not follow them in all things, and conform himself to the church of Rome, is judged to be a heretic. The inner sanctuary in the law was the station of priests, the place wherein they were when they should do sacrifice. And while he says the sanctuary must be cast out, he signifies figuratively that the Antichristian priests shall be thrown out. For the place is set for the thing contained therein. And where he says, cast out, he says it of those whom God has shut out; declare them to be cast out. For God does excommunicate; man pronounces and executes God’s judgment. The latter reading is of this sort: and the sanctuary which is without, cast out. So has the Spanish copy. But how shall you cast out that which was without before? Therefore I like, as I said, the former reading. But we reject this reading neither. For the hall that is without signifies the fellowship not communicating with the only altar, Christ, or with the true church of Christ, such as all this book shows the Popes to be with all their family. Moreover, the Pharisees and priests cast out the one who was born blind, John 9, that is to say, they excommunicated him for the confession of Christ, and the Lord says in John 15: “If anyone does not abide in me, he is cast out as a branch and withers.” Therefore, while John is commanded here to cast out the fellowship of priests, he is verily commanded to declare that those priests were excommunicated, which would be and seem the chief prelates of Christ’s church. He is also forbidden to meet this fellowship. For because God will not edify but destroy them, neither will have them honored among his. For he has rejected them. Who then will hereafter care much for the excommunication of those who are excommunicated? Wicked popes have excommunicated emperors, noblemen, and godly people; and, discharging their subjects of their fidelity, have set them in their princes’ places. The story of Gregory II against Leo Isauricus is known, and of Gregory VII against Henry IV, and of Innocent III against Frederick II, and of other bishops against right good princes. Doubtless, the chief string of the Popish tyranny has been excommunication, which the Lord here loosens.

The Lord does not conceal why he pronounces the priests, or inner choir, excommunicated; for it is given to the heathen. This phrasing is as forceful as saying that within the choir they do not act as priests or faithful ministers, but as gentiles who have taken possession of this place. But gentiles are rightly excluded from the fellowship of God and the church, where the Lord himself said in the Gospel: “If he does not hear the church, let him be to you as a heathen and publican.” Undoubtedly, those who are not in the Temple or church, or are in the inner choir—that is to say, those who are considered among the prelates of the church—yet do not hold to Christ, but are more conformed to the heathen than to Christians, are justly accounted excluded among the gentiles.

And now let us see why he identifies the Antichrist as the pope, with his followers and their images resembling the heathens. Those born of God hear the word of God and glorify it; those not yet born of God, but remain gentiles, not only do not hear God’s word, but also blaspheme it. So these men will not hear God’s word, and seek with all their endeavor how to frighten people away from the scriptures, which are God’s word. They say that they are obscure, doubtful, uncertain, and imperfect. Those who believe and cleave to them they call heretics, and the doctrine taken from them is heresy. Again, those who do not have Christ as their head, and as branches do not grow from the vine, have no communion with Christ and are gentiles. But such is the pope and his adherents, still persecuting Christ and all those who affirm Christ to be the only head of the church, Christ alone to be our righteousness and life, so that all the faithful are made fully complete by Christ. He who thus believes, they pronounce him a heretic.

Moreover, the gentiles worship idols, call upon creatures, suppose God to be honored with corruptible things such as gold, silver, and precious things, dedicated to the Temple and set up to beautify it. But what other thing do they do in the church at this day? You plainly see heathen temples when you see their churches. The life also of the gentiles is shameful and filthy; they are given to voluptuousness, full of surfeiting, addicted to filthy lust, they stink in whoredom, and exceed in gorgeous apparel and pampering of the body. See what things the Apostle writes of the life and conversation of the heathens in the fourth and fifth chapters to the Ephesians, and in the first chapter to the Romans, and in the first to the Corinthians, the fifth and sixth chapters.

\index{Resurrection}Now what the life of the pope is and of his spirituality, the thing itself openly testifies, that even for this cause alone they might and ought to be accounted among the excommunicated, the Apostle himself pronouncing the sentence of excommunication in the place which we have now cited, the first to the Corinthians, the fifth. We may add hereunto their Epicureanism. For if they value any religion, if they have any fear of God in them, why do they sell all things in the church – forgiveness of sins, heaven, Christ, the oblation of Christ, matrimony, ministry – briefly, all things? Why do they call into doubt diverse articles of our belief? What mean these doubtful disputations of the immortality of souls and the resurrection of bodies? Why do they make a mockery of the life everlasting?

Hereunto is added that they trample, yea spurn, the holy city; for this reason they may justly be declared excommunicated. This holy city is not that earthly Jerusalem, but the church of God, of which the holy city was a figure, as Paul explains in chapter 4 to the Galatians. For the earthly Jerusalem, according to the sayings of the prophets, having played its part, lies in ashes never to be restored. The Lord therefore signifies that the holy church of Christ should, through the tyranny of Antichrist and Antichristians, be trampled underfoot. And it signified more that he said to tread upon than if he had said to afflict and persecute. For treading is joined with the greatest contempt of him who is trodden on; and hereby is signified an extreme assault and wonderful cruelty of the enemies, which they practice on those they overcome, and have to use at their pleasure. We read in Daniel of the Romans: The beast had great iron teeth, eating and breaking small, and the rest treading under its feet. For wanton beasts are wont to tread with their feet such things as they cannot devour when they are full. And Solomon in the 27th of the Proverbs says, “A soul that is full treads the honeycomb.” Malachi in chapter 4, speaking of the joy of the godly, says, “They shall leap as calves of the herd, and they shall tread upon the wicked, which shall be as dust under the soles of your feet.” Briefly, John by treading signifies the oppression of the church joined with great tyranny, wantonness, and the exceeding great mockery and gladness of the wicked. And plainly seems to have alluded to these words of the godly prophet: “O God, the heathens have come into your inheritance, your holy temple they have defiled, and made Jerusalem an heap of stones. The dead bodies of your servants they have given to be meat unto the fowls of the air: and the flesh of your saints unto the beasts of the land. Their blood have they shed like water on every side of Jerusalem, and there was no man to bury them: and the rest that follows in the 78th Psalm.” And a little after in this chapter shall follow more things of the persecution of Antichrist. Neither shall these things be obscure, if you compare them with those which are done at this day in the church of Rome against the lovers of Christ’s gospel.

\index{Arethas of Caesarea}Here is shown a certain time in which the persecution of Antichrist would be cruel against the church, namely the space of two and forty months. Some, in accepting this, torment themselves marvelously. I suppose plainly that a certain time was assigned, and that not without cause, yet notwithstanding an uncertain time to be understood. A certain time therefore is assigned, that we might understand that God has appointed an end to their furies; which, as he himself alone knows, he would signify to his faithful, not in years, but in months only, for consolation. For we suffer more easily that which we perceive shall continue but a few months. This sense has also Arethas, after a sort, touched, writing thus: we suppose that the time of forty-two months expresses a shortening of time about the coming of Antichrist; for the affliction to be executed upon the lovers of God, Christ, very God, says that those days should be abbreviated. And these forty-two months are three years and a half, wherein it shall come to pass that the faithful, and the very tried, shall be trodden and suffer persecution. Thus says he.

\index{Apocalypse (Book of)}\index{Jerome, Saint}\index{Judgment, Last}Doubtless all commentators, in a manner being truly taught by this passage, have attributed to the kingdom of Antichrist and to his most cruel persecutions no more than three and a half years. For so many years make forty-two months, if one puts twelve months to a year. However, the Scripture and the thing itself speak that the kingdom of Antichrist should be a great deal longer. Whereupon I said that a certain time is assigned by the Apostle, and an uncertain time is understood: that is to say, all that same time that is reckoned from the fatal years 666, whereof is mentioned in the thirteenth chapter of the Apocalypse, until the last judgment. And why I expound a certain time by an uncertain, these are the causes. First, because the same number of months is put here in the thirteenth chapter and is ascribed to the old Roman Empire, so that in their tribulations the saints might understand and comfort themselves that there is an end appointed to their tyranny, which is known of God; and that the saints should no more be sorrowful than if they were constrained to abide their tyranny a few months only. Otherwise, if one should account from the first year of Julius Caesar and bring the course of time until that year wherein Odacer at Rome, all emperors of the west being taken away, was acknowledged as king, you shall not find only three years and a half, but about five hundred and seventeen years. If you shall bring the account from Julius to the empire taken away and given to the pope, you shall find about seven hundred and sixty-seven years. The later cause: for Daniel, the Lord Christ, and the Apostle Jerome agreeably do say that the persecution of Antichrist should last until the judgment. But who shall reckon to us the years and days of the last judgment? And therefore the number certain must be expounded by the uncertain, and must think that all things are numbered and prefixed in the counsel of God, which never neglects his faithful. To him be glory forevermore. Amen.
